\documentclass[border=10pt]{standalone}
\usepackage{fontspec}
\setmainfont{Noto Serif}[Ligatures=TeXOff]
\XeTeXgenerateactualtext=1
\newfontfamily\bengalifont{Noto Serif Bengali}[Script=Bengali, NFSSFamily=bnfam]
\usepackage[Latin,Bengali]{ucharclasses}
% group-free font switching (as in header-common.tex): robust inside braces, cells and quotes
\makeatletter
\def\BN@saved{\rmdefault}
\def\BN@bn{bnfam}
\newcommand\BNin{\edef\BN@tmp{\f@family}\ifx\BN@tmp\BN@bn\else\let\BN@saved\BN@tmp\fi\fontfamily{bnfam}\selectfont}
\newcommand\BNout{\fontfamily{\BN@saved}\selectfont}
\makeatother
\setTransitionsFor{Bengali}{\BNin}{\BNout}
% Typographic quotes and dashes are not in a ucharclasses class (class 0), so a Bangla word right
% after “ or ‘ would not trigger the Bangla font. Treat them as Latin so that every switch fires.
\XeTeXcharclass"2018=\XeTeXcharclass`A
\XeTeXcharclass"2019=\XeTeXcharclass`A
\XeTeXcharclass"201C=\XeTeXcharclass`A
\XeTeXcharclass"201D=\XeTeXcharclass`A
\XeTeXcharclass"2013=\XeTeXcharclass`A
\XeTeXcharclass"2014=\XeTeXcharclass`A
\XeTeXcharclass"2026=\XeTeXcharclass`A
\usepackage{newunicodechar}
\newunicodechar{।}{{\bengalifont\char"0964}}
\newunicodechar{॥}{{\bengalifont\char"0965}}

\usepackage[table]{xcolor}
\usepackage{booktabs,tabularx,array}
\definecolor{bapilrule}{HTML}{1F2937}
\definecolor{bapilstripe}{HTML}{F3F4F6}
\renewcommand{\arraystretch}{1.28}
\setlength{\tabcolsep}{5pt}

\begin{document}
\begin{minipage}{17cm}
{\large\bfseries সারণি ২৪।}\ {\large নথির মেটাডেটা ক্ষেত্র (policy\_document স্কিমা)}\par\vspace{7pt}
\small
\rowcolors{2}{bapilstripe}{white}
\begin{tabularx}{\linewidth}{@{}>{\hsize=0.762\hsize\raggedright\arraybackslash}X>{\hsize=1.667\hsize\raggedright\arraybackslash}X>{\hsize=0.429\hsize\raggedright\arraybackslash}X>{\hsize=1.143\hsize\raggedright\arraybackslash}X@{}}
\arrayrulecolor{bapilrule}\toprule
\rowcolor{white}\textbf{ক্ষেত্র} & \textbf{বিবরণ} & \textbf{আবশ্যিক} & \textbf{উদাহরণ} \\
\midrule
id & কর্পাসের ভেতরে স্থায়ী শনাক্তকারী & হ্যাঁ & bd-act-2009-20 \\
title & প্রাপ্তিসাপেক্ষে বাংলা ও ইংরেজি শিরোনাম & হ্যাঁ & তথ্য অধিকার আইন, ২০০৯ \\
publisher & নথি প্রকাশকারী প্রতিষ্ঠান & হ্যাঁ & লেজিসলেটিভ ও সংসদ বিষয়ক বিভাগ \\
issuing\_body & দলিল জারিকারী সংস্থা & না & জাতীয় সংসদ \\
document\_type & আইন, অধ্যাদেশ, বিধি, গেজেট, বাজেট, পরিকল্পনা, পরিসংখ্যান, সংসদীয় কাগজ, নীতি, প্রতিবেদন & হ্যাঁ & act \\
date\_issued & জারির তারিখ & হ্যাঁ & 2009-04-06 \\
version & সংস্করণ বা সংশোধনের অবস্থা & না & ২০২৬ পর্যন্ত সংশোধিত \\
language & টেক্সটের ভাষা & হ্যাঁ & bn, en \\
source\_url & মূল নথি যেখান থেকে সংগৃহীত & হ্যাঁ & http:/\allowbreak{}/\allowbreak{}bdlaws.minlaw.gov.bd/\allowbreak{}... \\
retrieved & সংগ্রহের তারিখ & হ্যাঁ & 2027-02-01 \\
reuse\_terms & লাইসেন্স বা পুনর্ব্যবহারের শর্ত & না & সরকারি প্রকাশনা; শর্ত নিশ্চিত \\
data\_class & DC0, DC1 বা DC2 & হ্যাঁ & DC0 \\
sha256 & মূল ফাইলের চেকসাম & হ্যাঁ & ৬৪টি হেক্সাডেসিমাল অক্ষর \\
extraction & পদ্ধতি (জন্মগতভাবে ডিজিটাল বা OCR), মানের স্কোর, সরঞ্জাম & হ্যাঁ & ocr, 0.93, tesseract 5 \\
supersedes & আগের সংস্করণের শনাক্তকারী & না & bd-act-2009-20-v1 \\
\bottomrule
\end{tabularx}
\par\vspace{5pt}{\footnotesize\itshape উদাহরণগুলো দৃষ্টান্তমূলক। যন্ত্র-পাঠযোগ্য স্কিমা: implementation/\allowbreak{}schemas/\allowbreak{}policy\_document.schema.json।}
\par\vspace{3pt}{\scriptsize\color{gray} উৎস: BAPIL টেকনিক্যাল ব্লুপ্রিন্ট, সংস্করণ ১.০ (২০২৬)। CC BY 4.0।}
\end{minipage}
\end{document}
